\chapter{Generating Functionals and Ward Identities}

In quantum field theory, the path integral formalism provides a unified framework for encoding the complete quantum dynamics of a system into functional objects known as generating functionals. By coupling the quantum fields to external sources, one can systematically generate arbitrary time-ordered correlation functions through functional differentiation. Furthermore, this functional language allows one to isolate the physically relevant connected and one-particle-irreducible (1PI) diagrams, while providing a powerful algebraic arena to translate continuous symmetries of the classical action into exact quantum relations, known as Ward--Takahashi or Slavnov--Taylor identities.

\section{Generating Functionals}

Let us systematically describe the fundamental generating functionals obtained by quantizing a generic classical action \(S[\phi]\):
\begin{equation} \label{eq:partition_function_Z}
    Z[J] = \int \mathcal{D}\phi \, e^{iS[\phi] + i \int d^4x \, J(x)\phi(x)} \equiv e^{iW[J]} \, ,
\end{equation}
\begin{equation} \label{eq:connected_generator_W}
    W[J] = -i \ln Z[J] \, ,
\end{equation}
\begin{equation} \label{eq:effective_action_Gamma}
    \Gamma[\varphi] = \underset{J}{\text{ext}} \left\{ W[J] - \int d^4x \, J(x)\varphi(x) \right\} \, .
\end{equation}

Each of these three functionals plays a distinct and crucial structural role in the perturbative expansion:
\begin{itemize}
    \item \(Z[J]\) is the generating functional of full \textit{correlation functions} (including disconnected vacuum and multi-particle processes),
    \item \(W[J]\) is the generating functional of \textit{connected correlation functions},
    \item \(\Gamma[\varphi]\), defined via the Legendre transform of \(W[J]\) with respect to the classical background field \(\varphi(x)\), is the generating functional of 1PI (\textit{one-particle-irreducible}) functions, also universally known as the \textit{quantum effective action}. By 1PI functions, one specifically means those correlation functions given diagrammatically by graphs that cannot be disconnected into two separate pieces by cutting any single internal line, and which have their external lines amputated (i.e., without external propagators associated with them).
\end{itemize}

As an explicit prototype to test and illustrate the general discussion, one may keep in mind the classical action of a real scalar field with a generic polynomial interaction:
\begin{equation} \label{eq:scalar_polynomial_action}
    S[\phi] = \int d^4x \Bigg[ \underbrace{-\frac{1}{2}\partial_\mu\phi\partial^\mu\phi - \frac{1}{2}m^2\phi^2}_{\text{free action}} - \sum_{n \ge 3} \underbrace{\frac{g_n}{n!}\phi^n}_{\text{classical vertices}} - g_1\phi - g_0 \Bigg] \, ,
\end{equation}
where \(g_n\) are the various interaction coupling constants. We have explicitly included here also the linear coupling \(g_1\) to account for the possible existence of a one-point coupling constant (if present, it signals physically that the configuration \(\phi = 0\) is not a true minimum of the classical potential; however, it can always be eliminated by suitably shifting the field \(\phi\)) and the constant energy density \(g_0\) that may become physically relevant when coupling the theory to gravity (as it directly contributes to the cosmological constant). One could also naturally introduce derivative interactions to further enlarge the space of couplings, as famously exemplified in the nonlinear sigma model or by the cubic self-interactions of a non-abelian Yang--Mills gauge field.

Correlation functions contain the essential dynamical information needed to extract physical observables in Quantum Field Theory. A primary example is given by a scattering cross section, which is constructed directly from on-shell scattering amplitudes that, in turn, are systematically obtained from time-ordered correlation functions through the LSZ reduction formula.

\subsection{The Classical Action}

The quantum correlation functions generated by the functional path integral in \eqref{eq:partition_function_Z} are evaluated perturbatively using Feynman diagrams built from the free propagator and the elementary vertices derived directly from the (eventually gauge-fixed) classical action, such as the scalar prototype in \eqref{eq:scalar_polynomial_action}. To systematize the structural information contained within the classical action, one formally defines the \textit{classical vertices} as the Taylor coefficients of the functional Taylor expansion of the classical action around the vanishing field configuration.

Using DeWitt's hypercondensed notation (where a single composite index \(i\) encompasses both discrete internal indices and continuous spacetime coordinates, so that index summation implicitly includes spacetime integration, as in \(J_i \phi^i = \int d^4x \, J(x)\phi(x)\)), we define the classical vertices \(V^{(n)}\) as the functional derivatives of the classical action multiplied by \(i\):
\begin{equation} \label{eq:classical_vertices_def}
    \boxed{V^{(n)}_{i_1 \cdots i_n} = \left. \frac{\delta^n iS[\phi]}{\delta\phi^{i_1} \cdots \delta\phi^{i_n}} \right|_{\phi=0}} \, ,
\end{equation}
which immediately allows us to write the classical action as a formal functional Taylor expansion:
\begin{equation} \label{eq:action_taylor_expansion}
    iS[\phi] = \sum_{n=0}^\infty \frac{1}{n!} \phi^{i_1} \cdots \phi^{i_n} V^{(n)}_{i_1 \cdots i_n} \, ,
\end{equation}
where the imaginary unit \(i\) is inserted by convention because it is precisely the combination \(iS[\phi]\) that enters the phase of the path integral weight \(e^{iS[\phi]}\). Strictly speaking, the terminology ``classical vertices'' applies properly only for \(n \ge 3\), whereas the \(n=2\) term defines the inverse free propagator (and the \(n=0, 1\) terms represent the vacuum energy and lower-point tadpole vertices, when present).

For instance, returning to the concrete scalar example of \eqref{eq:scalar_polynomial_action}, one gets for the cubic interaction vertex (\(n=3\)) in coordinate space:
\begin{equation} \label{eq:cubic_vertex_coord}
    V^{(3)}_{i_1 i_2 i_3} \equiv V^{(3)}(x_1, x_2, x_3) = -ig_3 \int d^4x \, \delta^4(x - x_1)\delta^4(x - x_2)\delta^4(x - x_3) \, .
\end{equation}
This expression reflects a strictly \textit{local vertex}, as only Dirac delta functions enforcing the coincidence of spacetime points appear. In more general theories, a finite number of spacetime derivatives acting on the delta functions are also allowed for local interactions. Taking the Fourier transform of a coordinate-space vertex directly gives the corresponding Feynman rule in momentum space:
\begin{equation} \label{eq:cubic_vertex_momentum}
    \begin{aligned}
        V^{(3)}(k_1, k_2, k_3) & = \int d^4x_1 d^4x_2 d^4x_3 \, e^{ik_1 \cdot x_1} e^{ik_2 \cdot x_2} e^{ik_3 \cdot x_3} V^{(3)}(x_1, x_2, x_3) \\
                               & = -ig_3 \, (2\pi)^4 \delta^4(k_1 + k_2 + k_3) \, .
    \end{aligned}
\end{equation}
Up to the overall factor of \((2\pi)^4 \delta^4(k_1 + k_2 + k_3)\) ensuring total four-momentum conservation at the interaction point, this yields precisely the algebraic factor \(-ig_3\) associated with the standard Feynman rule for the cubic vertex.

\TODO{add image of three point vertex with incoming momenta k1 k2 k3 and factor minus i g3}

\subsection{The Partition Function: Generator of Correlation Functions}

Let us now formally describe the full quantum correlation functions. They are generated by taking functional derivatives of the partition function \(Z[J]\) with respect to the external source \(J\). Here, for algebraic simplicity in tracking disconnected pieces, it is convenient to initially define the correlation functions \textit{without} dividing by the vacuum-to-vacuum normalization factor \(Z[0]\):
\begin{equation} \label{eq:unnormalized_correlators}
    \boxed{\langle \phi^{i_1} \cdots \phi^{i_n} \rangle = \left. \left(\frac{1}{i}\right)^n \frac{\delta^n Z[J]}{\delta J_{i_1} \cdots \delta J_{i_n}} \right|_{J=0}} \, .
\end{equation}

\TODO{add image of vacuum bubbles expansion for Z at zero source}

In particular, setting \(n=0\), the quantity \(Z[0] = \langle 1 \rangle\) gives the overall vacuum normalization of the path integral. Diagrammatically, it can be identified as the sum of all vacuum bubble graphs (diagrams with no external legs), encompassing both connected and disconnected topologies, where the combinatoric weight of each graph is not explicitly reported in the sketch, as it can be systematically obtained by computing the number of equivalent Wick contractions.

When evaluated for \(n \ge 1\), these full correlation functions \(\langle \phi^{i_1} \cdots \phi^{i_n} \rangle\) inevitably contain unphysical disconnected pieces, such as vacuum bubbles multiplying the physical process or sub-processes occurring independently at large spacetime separations. It is therefore physically necessary and mathematically convenient to systematically eliminate them and isolate the purely \textit{connected correlation functions}.

\subsection{Generator of Connected Correlation Functions}

Connected correlation functions are directly generated by taking functional derivatives of the functional \(W[J] = -i \ln Z[J]\) as follows:
\begin{equation} \label{eq:connected_correlators_def}
    \boxed{\langle \phi^{i_1} \cdots \phi^{i_n} \rangle_c = \left. \left(\frac{1}{i}\right)^n \frac{\delta^n iW[J]}{\delta J_{i_1} \cdots \delta J_{i_n}} \right|_{J=0}} \, .
\end{equation}

\subsubsection{Connected One-Point Function}

To concretely understand why \(W[J]\) isolates connected topologies, let us first look at the connected one-point function (the vacuum expectation value of the field):
\begin{equation} \label{eq:connected_one_point}
    \langle \phi^i \rangle_c = \left. \frac{\delta W[J]}{\delta J_i} \right|_{J=0} = \left. \frac{\delta}{\delta J_i} \left(\frac{1}{i} \ln Z[J]\right) \right|_{J=0} = \left. \frac{1}{Z[0]} \left(\frac{1}{i} \frac{\delta Z[J]}{\delta J_i}\right) \right|_{J=0} = \frac{\langle \phi^i \rangle}{\langle 1 \rangle} \, ,
\end{equation}
where we immediately see that the natural appearance of the division by \(\langle 1 \rangle = Z[0]\) precisely factors out and removes all the disconnected vacuum bubble pieces that otherwise multiply the diagram. We can algebraically rewrite this relation also as:
\begin{equation} \label{eq:one_point_factorization}
    \langle \phi^i \rangle = \langle \phi^i \rangle_c \, \langle 1 \rangle \, ,
\end{equation}
which shows how the full unnormalized one-point function factors into the connected one-point tadpole times the sum of all vacuum bubbles. Diagrammatically, we depict this factorization as:

\TODO{add image of diagrammatic decomposition of the one point function into connected part and vacuum bubbles}

\subsubsection{Connected Two-Point Function}

Let us now proceed to look at the connected 2-point function by taking a second functional derivative:
\begin{equation} \label{eq:connected_two_point}
    \begin{aligned}
        \langle \phi^i \phi^j \rangle_c & = \left. \left(\frac{1}{i}\right)^2 \frac{\delta^2 iW[J]}{\delta J_i \delta J_j} \right|_{J=0} = \left. \left(\frac{1}{i}\frac{\delta}{\delta J_i}\right) \left(\frac{1}{i}\frac{\delta}{\delta J_j}\right) \ln Z[J] \right|_{J=0}                                               \\
                                        & = \left. \left( \frac{1}{Z[J]} \frac{1}{i^2} \frac{\delta^2 Z[J]}{\delta J_i \delta J_j} - \frac{1}{Z^2[J]} \frac{1}{i}\frac{\delta Z[J]}{\delta J_i} \frac{1}{i}\frac{\delta Z[J]}{\delta J_j} \right) \right|_{J=0}                                                            \\
                                        & = \frac{\langle \phi^i \phi^j \rangle}{\langle 1 \rangle} - \frac{\langle \phi^i \rangle}{\langle 1 \rangle} \frac{\langle \phi^j \rangle}{\langle 1 \rangle} = \frac{\langle \phi^i \phi^j \rangle}{\langle 1 \rangle} - \langle \phi^i \rangle_c \langle \phi^j \rangle_c \, .
    \end{aligned}
\end{equation}
We clearly see from this result that the first term on the right-hand side is the normalized 2-point function stripped of all disconnected vacuum bubbles, while the second term explicitly subtracts the disconnected 2-point contributions that factorize into the product of two independent connected 1-point functions.

\TODO{add image of diagrammatic representation of the connected two point function}

One can invert and rewrite this relation in terms of the full unnormalized correlator as:
\begin{equation} \label{eq:two_point_decomposition}
    \langle \phi^i \phi^j \rangle = \langle 1 \rangle \big[ \langle \phi^i \phi^j \rangle_c + \langle \phi^i \rangle_c \langle \phi^j \rangle_c \big] \, .
\end{equation}

\subsubsection{Connected Three-Point Function}

Let us continue iteratively and consider the connected 3-point function by applying a third functional derivative:
\begin{equation} \label{eq:connected_three_point_derivation}
    \begin{aligned}
        \langle \phi^1 \phi^2 \phi^3 \rangle_c & = \left. \left(\frac{1}{i}\right)^3 \frac{\delta^3 iW[J]}{\delta J_1 \delta J_2 \delta J_3} \right|_{J=0} = \left. \left(\frac{1}{i}\frac{\delta}{\delta J_1}\right) \left(\frac{1}{i^2}\frac{\delta^2}{\delta J_2 \delta J_3}\right) \ln Z[J] \right|_{J=0}            \\
                                               & = \left. \left(\frac{1}{i}\frac{\delta}{\delta J_1}\right) \left( \frac{1}{Z[J]} \frac{1}{i^2} \frac{\delta^2 Z[J]}{\delta J_2 \delta J_3} - \frac{1}{Z^2[J]} \frac{1}{i}\frac{\delta Z[J]}{\delta J_2} \frac{1}{i}\frac{\delta Z[J]}{\delta J_3} \right) \right|_{J=0} \\
                                               & = \Bigg[ \frac{1}{Z[J]} \frac{1}{i^3} \frac{\delta^3 Z[J]}{\delta J_1 \delta J_2 \delta J_3} - \left( \frac{1}{Z^2[J]} \frac{1}{i}\frac{\delta Z[J]}{\delta J_1} \frac{1}{i^2}\frac{\delta^2 Z[J]}{\delta J_2 \delta J_3} + \text{cyclic perm. of } 1, 2, 3 \right)     \\
                                               & \quad + 2 \frac{1}{Z^3[J]} \frac{1}{i^3} \frac{\delta Z[J]}{\delta J_1} \frac{\delta Z[J]}{\delta J_2} \frac{\delta Z[J]}{\delta J_3} \Bigg] \Bigg|_{J=0}                                                                                                               \\
                                               & = \Bigg[ \frac{1}{Z[J]} \frac{1}{i^3} \frac{\delta^3 Z[J]}{\delta J_1 \delta J_2 \delta J_3} - \left( \frac{\delta W[J]}{\delta J_1} \frac{1}{i^2} \frac{\delta^2 iW[J]}{\delta J_2 \delta J_3} + \text{cyclic perm. of } 1, 2, 3 \right)                               \\
                                               & \quad - \frac{\delta W[J]}{\delta J_1} \frac{\delta W[J]}{\delta J_2} \frac{\delta W[J]}{\delta J_3} \Bigg] \Bigg|_{J=0} \, ,
    \end{aligned}
\end{equation}
where in obtaining the final step, we systematically substituted the second line of \eqref{eq:connected_two_point} as well as \eqref{eq:connected_one_point} to re-express the intermediate derivatives in terms of \(W[J]\). Evaluating at \(J=0\), this algebraic relationship represents:
\begin{equation} \label{eq:connected_three_point}
    \langle \phi^1 \phi^2 \phi^3 \rangle_c = \frac{\langle \phi^1 \phi^2 \phi^3 \rangle}{\langle 1 \rangle} - \Big( \langle \phi^1 \rangle_c \langle \phi^2 \phi^3 \rangle_c + \text{cyclic perm. of } 1, 2, 3 \Big) - \langle \phi^1 \rangle_c \langle \phi^2 \rangle_c \langle \phi^3 \rangle_c \, .
\end{equation}

\TODO{add image of diagrammatic representation of the connected three point function}

Equivalently, one may invert this expression and write it as:
\begin{equation} \label{eq:three_point_decomposition}
    \langle \phi^1 \phi^2 \phi^3 \rangle = \langle 1 \rangle \Big[ \langle \phi^1 \phi^2 \phi^3 \rangle_c + \langle \phi^1 \phi^2 \rangle_c \langle \phi^3 \rangle_c + \langle \phi^1 \phi^3 \rangle_c \langle \phi^2 \rangle_c + \langle \phi^2 \phi^3 \rangle_c \langle \phi^1 \rangle_c + \langle \phi^1 \rangle_c \langle \phi^2 \rangle_c \langle \phi^3 \rangle_c \Big] \, ,
\end{equation}
which transparently describes how the full 3-point correlation function is combinatorially obtained by putting together all possible partitions into connected 3-point and lower-point correlation functions, multiplied overall by the vacuum bubble normalization \(\langle 1 \rangle\).

The above cluster decomposition relations generalize systematically to all higher-point functions:
\begin{equation} \label{eq:general_connected_correlator}
    \langle \phi^1 \phi^2 \cdots \phi^n \rangle_c = \frac{\langle \phi^1 \phi^2 \cdots \phi^n \rangle}{\langle 1 \rangle} - (\text{all possible disconnected } n\text{-point functions}) \, .
\end{equation}

\subsection{The Effective Action: Generator of 1PI Correlation Functions}

Let us now turn to discuss the quantum effective action \(\Gamma[\varphi]\) in detail and convince ourselves that it systematically generates one-particle-irreducible (1PI) correlation functions. As previously introduced, the latter are those correlation functions identified diagrammatically by graphs that cannot be disconnected into two separate pieces by cutting a single internal line, and which have their external lines fully amputated (i.e., without external propagators associated with them).

Mathematically, the effective action is defined by the functional Legendre transform of the connected generating functional \(W[J]\):
\begin{equation} \label{eq:legendre_transform_gamma}
    \Gamma[\varphi] = \text{ext}_J \left[ W[J] - \int d^4x \, J(x)\varphi(x) \right] \, .
\end{equation}
The stationarity (extremization) condition means that one must evaluate the right-hand side at its stationary point with respect to the source \(J\), which is determined by the functional equation:
\begin{equation} \label{eq:stationarity_condition_J}
    \frac{\delta W[J]}{\delta J(x)} - \varphi(x) = 0 \, .
\end{equation}
This relation produces a functional \(\varphi = \varphi[J]\), which in principle should be inverted to find the source as a functional of the field, \(J = J[\varphi]\). Substituting this inverted relation back into \eqref{eq:legendre_transform_gamma} then expresses \(\Gamma[\varphi]\) strictly as a functional of the classical background field \(\varphi\) only:
\begin{equation} \label{eq:gamma_explicit_phi}
    \Gamma[\varphi] = W[J[\varphi]] - \int d^4x \, J[\varphi](x) \, \varphi(x) \, .
\end{equation}

The Legendre transform structure also directly implies the conjugate functional relation:
\begin{equation} \label{eq:conjugate_legendre_relation}
    \frac{\delta \Gamma[\varphi]}{\delta \varphi(x)} = -J(x) \, ,
\end{equation}
which is readily verified by a direct functional chain-rule calculation:
\begin{equation} \label{eq:legendre_verification}
    \begin{aligned}
        \frac{\delta \Gamma[\varphi]}{\delta \varphi(x)} & = \frac{\delta}{\delta \varphi(x)} \left[ W[J] - \int d^4y \, J(y)\varphi(y) \right]                                                                                                                                 \\
                                                         & = \int d^4y \, \frac{\delta W[J]}{\delta J(y)} \frac{\delta J(y)}{\delta \varphi(x)} - \int d^4y \, \frac{\delta J(y)}{\delta \varphi(x)} \varphi(y) - \int d^4y \, J(y) \frac{\delta \varphi(y)}{\delta \varphi(x)} \\
                                                         & = - \int d^4y \, J(y) \delta^4(y - x) = -J(x) \, ,
    \end{aligned}
\end{equation}
where the stationarity condition of \eqref{eq:stationarity_condition_J} has been used to cancel out and eliminate the first two terms of the second line.

Thus, we have established the fundamental dual properties of the Legendre transformation:
\begin{equation} \label{eq:basic_legendre_properties}
    \begin{aligned}
        \frac{\delta W[J]}{\delta J(x)}                  & = \varphi(x) \, , \\
        \frac{\delta \Gamma[\varphi]}{\delta \varphi(x)} & = -J(x) \, ,
    \end{aligned}
\end{equation}
which give \(\varphi(x)\) as a functional of \(J(x)\) as well as \(J(x)\) as a functional of \(\varphi(x)\). These reciprocal equations form the algebraic backbone that will be used to find the main diagrammatic properties of the effective action.

\subsubsection{The Mean Field and One-Point Function}

Before starting the higher-order derivative analysis, let us take a closer look at the physical definition of the \textit{mean field} (which is another standard name for the background variable \(\varphi\)):
\begin{equation} \label{eq:mean_field_def}
    \varphi(x) = \frac{\delta W[J]}{\delta J(x)} = \langle \phi(x) \rangle_{c, J} \, ,
\end{equation}
which represents the connected one-point function evaluated in the presence of a nonzero external source \(J(x)\). It is a spacetime function that depends functionally on the source \(J(x)\). If we turn off the external source by setting \(J(x) = 0\), we recover the standard connected one-point function (the vacuum expectation value, or VEV) of the quantum theory:
\begin{equation} \label{eq:vev_definition}
    \left. \varphi(x) \right|_{J=0} = \langle \phi(x) \rangle_c \equiv v \, ,
\end{equation}
which reduces to a spacetime-independent constant \(v\) in a translationally and Lorentz invariant vacuum state. For a scalar field theory, a nonzero vacuum expectation value \(v \neq 0\) may physically signal the spontaneous breaking of an internal symmetry.

\subsubsection{1PI Two-Point Function}

Now, let us consider the 2-point function generated by \(\Gamma[\varphi]\). Using DeWitt's hypercondensed notation, we differentiate the conjugate Legendre relation a second time:
\begin{equation} \label{eq:gamma_two_point_initial}
    \frac{\delta^2 \Gamma[\varphi]}{\delta \varphi^1 \delta \varphi^2} = \frac{\delta}{\delta \varphi^1} \frac{\delta \Gamma[\varphi]}{\delta \varphi^2} = \frac{\delta}{\delta \varphi^1} (-J_2) = -\frac{\delta J_2}{\delta \varphi^1} \, .
\end{equation}
The functional Jacobians \(\delta \varphi / \delta J\) and \(\delta J / \delta \varphi\) are inverse integral kernels of one another. To see this explicitly, note that functionally differentiating \(\varphi^j = \delta W[J] / \delta J_j\) with respect to the source gives:
\begin{equation} \label{eq:jacobian_w_two_point}
    \frac{\delta \varphi^j}{\delta J_k} = \frac{\delta^2 W[J]}{\delta J_j \delta J_k} \, .
\end{equation}
Multiplying the preceding expression for the second derivative of \(\Gamma\) by this Jacobian and applying the functional chain rule, we obtain the fundamental identity:
\begin{equation} \label{eq:gamma_w_inverse_condensed}
    \frac{\delta^2 \Gamma}{\delta \varphi^i \delta \varphi^j} \frac{\delta^2 W}{\delta J_j \delta J_k} = -\frac{\delta J_i}{\delta \varphi^j} \frac{\delta \varphi^j}{\delta J_k} = -\delta_i^k \, .
\end{equation}

Note that the bosonic second-derivative kernels are manifestly symmetric in their two indices. In explicit spacetime coordinate notation, the contraction over the repeated composite index \(j\) translates into a spacetime integral:
\begin{equation} \label{eq:gamma_w_inverse_spacetime}
    \int d^4y \, \frac{\delta^2 \Gamma}{\delta \varphi(x) \delta \varphi(y)} \frac{\delta^2 W}{\delta J(y) \delta J(z)} = -\delta^4(x - z) \, .
\end{equation}
Thus, the inverse is understood as an inverse integral kernel (operator inverse), rather than a naive pointwise reciprocal. Consequently, the second functional derivative of \(\Gamma\) is precisely minus the inverse of the connected propagator:
\begin{equation} \label{eq:exact_propagator_inversion}
    \boxed{\frac{\delta^2 \Gamma[\varphi]}{\delta \varphi^i \delta \varphi^j} \frac{\delta^2 W[J]}{\delta J_j \delta J_k} = -\delta_i^k} \, ,
\end{equation}
so that \(\Gamma^{(2)}\) represents the \textit{exact quantum kinetic term} (incorporating all self-energy loop corrections), whose operator inverse defines the exact connected propagator \(W^{(2)}\).

\subsubsection{1PI Three-Point Function}

The structural properties of the higher \(n\)-point functions generated by \(\Gamma[\varphi]\) follow systematically by taking further functional derivatives of the inversion relation in \eqref{eq:exact_propagator_inversion}. For instance, let us compute the \(\frac{\delta}{\delta \varphi^1}\) derivative of this equation, using simple numbers \((1, 2, 3)\) for the external composite indices to make the resulting tensor equations more readable:
\begin{equation} \label{eq:derivative_inversion_relation}
    \frac{\delta}{\delta \varphi^1} \left( \frac{\delta^2 \Gamma[\varphi]}{\delta \varphi^2 \delta \varphi^i} \frac{\delta^2 W[J]}{\delta J_i \delta J_3} \right) = 0 \, .
\end{equation}
Using the product and functional chain rules (recalling that \(\frac{\delta}{\delta \varphi^1} = \frac{\delta J_j}{\delta \varphi^1} \frac{\delta}{\delta J_j}\)), one finds:
\begin{equation} \label{eq:three_point_chain_rule}
    \frac{\delta^3 \Gamma[\varphi]}{\delta \varphi^1 \delta \varphi^2 \delta \varphi^i} \frac{\delta^2 W[J]}{\delta J_i \delta J_3} + \frac{\delta^2 \Gamma[\varphi]}{\delta \varphi^2 \delta \varphi^i} \frac{\delta^3 W[J]}{\delta J_j \delta J_i \delta J_3} \frac{\delta J_j}{\delta \varphi^1} = 0 \, ,
\end{equation}
which, after contracting with two exact propagators \(\frac{\delta^2 W}{\delta J \delta J}\) and using \eqref{eq:exact_propagator_inversion} to isolate the third derivative of \(W[J]\), can be rewritten as:
\begin{equation} \label{eq:connected_three_point_1pi}
    \frac{\delta^3 W[J]}{\delta J_1 \delta J_2 \delta J_3} = \frac{\delta^3 \Gamma[\varphi]}{\delta \varphi^i \delta \varphi^j \delta \varphi^k} \frac{\delta^2 W[J]}{\delta J_i \delta J_1} \frac{\delta^2 W[J]}{\delta J_j \delta J_2} \frac{\delta^2 W[J]}{\delta J_k \delta J_3} \, .
\end{equation}

\TODO{add image of diagrammatic decomposition of the connected three point function into a 1PI vertex and three exact propagators}

This relation, depicted diagrammatically in the figure above, demonstrates that the third derivative \(\Gamma^{(3)}\) is an irreducible vertex, i.e., it is strictly 1PI. In this diagrammatic structure, there is a central 1PI irreducible vertex, which is fully amputated and cannot be disconnected by cutting any single internal line. The 1PI vertex is connected to the external legs \((1, 2, 3)\) by three full (exact) propagators \(W^{(2)}\). If we cut a single line connecting the 1PI vertex to the external legs, we can disconnect an external point from the diagram, but we have already explicitly accounted for the full propagators connecting the central 1PI core to the external legs. Because the external lines of \(\Gamma^{(3)}\) itself do not contain the propagators, cutting any of its internal lines will not break the vertex diagram into smaller pieces, i.e., it remains connected. The 1PI vertex is thus the fundamental building block and generator of all 1PI correlation functions.

\subsubsection{1PI Four-Point Function}

Proceeding iteratively to the 4-point function, we compute it by acting with the functional derivative \(\frac{\delta}{\delta J_4}\) directly on \eqref{eq:connected_three_point_1pi} and get, in an obvious condensed notation where superscripts/subscripts denote functional derivatives:
\[
    \begin{aligned}
        W^{1234} & = \Gamma_{ijkm} W^{i1} W^{j2} W^{k3} W^{m4} + \Big( \Gamma_{ijk} W^{i14} W^{j2} W^{k3} + \text{cyclic perm of } 1, 2, 3 \Big)                                               \\
                 & = \Gamma_{ijkm} W^{i1} W^{j2} W^{k3} W^{m4} + \Big[ \Gamma_{ijk} \left( W^{mi} \Gamma_{mnp} W^{n1} W^{p4} \right) W^{j2} W^{k3} + \text{cyclic perm of } 1, 2, 3 \Big] \, ,
    \end{aligned}
\]
where in passing to the second line we have substituted \eqref{eq:connected_three_point_1pi} once again to expand the three-point function \(W^{i14}\). This algebraic relation is much better interpreted in a graphical form, as shown in the figure below.

\TODO{add image of diagrammatic decomposition of the connected four point function into 1PI four point vertex and exchange tree diagrams}

\textit{Note on the diagrammatic representation:} In the original handwritten sketch of the connected 4-point function, care must be taken with the final exchange diagram (the last cyclic permutation at the bottom right): an exact propagator \(-iW^{(2)}\) must be present on every external leg, including leg 4 (which was inadvertently drawn directly attached to the \(i\Gamma^{(3)}\) vertex without its external propagator blob).

As the equation and diagram clearly reveal, the full connected 4-point function is obtained by summing tree-level topologies built entirely from proper 1PI vertices (\(\Gamma^{(4)}\) and pairs of \(\Gamma^{(3)}\)) connected by exact propagators \(W^{(2)}\). This remarkable construction generalizes to arbitrary connected \(n\)-point functions.

Thus, the quantum effective action \(\Gamma[\varphi]\) is used mathematically just like a classical action: its proper vertices and exact propagator already incorporate all quantum loop effects, so full connected correlation functions are completely reconstructed from tree diagrams alone, without adding any further loops. It is called the \textit{effective action} precisely because one does not have to add other loops or topological structures to describe quantum effects: all virtual radiative corrections are already fully resummed and incorporated into this single generating functional. The resulting effective vertices are then obtained simply by taking functional derivatives of \(\Gamma[\varphi]\), in complete analogy to how classical vertices are extracted from the classical action \(S[\phi]\).

To summarize, the proper vertices (i.e., the amputated 1PI functions) are generated by functionally differentiating the effective action:
\begin{equation} \label{eq:proper_vertices_summary}
    \boxed{\Gamma^{(n)}_{i_1 \cdots i_n} := \left. \frac{\delta^n \Gamma[\varphi]}{\delta \varphi^{i_1} \cdots \delta \varphi^{i_n}} \right|_{\varphi = v}} \, .
\end{equation}
The corresponding Minkowski-space Feynman-rule factor for a proper vertex is \(i\Gamma^{(n)}\), which should be directly compared with the bare classical vertex definition \(V^{(n)} = \left. \frac{\delta^n iS[\phi]}{\delta \phi^{i_1} \cdots \delta \phi^{i_n}} \right|_{\phi=0}\) introduced earlier. Here \(v\) denotes the vacuum expectation value of the field; after appropriately shifting the quantum field so that \(v = 0\), the functional derivatives are evaluated simply at \(\varphi = 0\). Unlike the local classical vertices \(V^{(n)}\), the proper quantum vertices \(\Gamma^{(n)}\) are generally nonlocal functionals of spacetime and contain all quantum effects, including arbitrary virtual-particle loops.

In general, the loop integrals entering the effective action exhibit ultraviolet divergences and must be made finite by applying the renormalization procedure. Once \(\Gamma[\varphi]\) is renormalized, one can directly construct the finite connected correlation functions and proceed to study physical processes such as on-shell scattering amplitudes.

\section{Path Integrals, Equations of Motion, Ward Identities}

Let us now comment on the fundamental role of the classical equations of motion and symmetries within the quantum context using the path integral formalism. Throughout this section, we employ the Euclidean form of the path integral (obtained via a Wick rotation to imaginary time), where the oscillatory weight \(e^{iS}\) is replaced by the exponentially damped weight \(e^{-S}\).

The partition function in the presence of external sources may be written equivalently as:
\begin{equation} \label{eq:path_integral_mute_variable}
    Z[J] = \int \mathcal{D}\phi \, e^{-S[\phi] + \int d^4x \, J(x)\phi(x)} = \int \mathcal{D}\phi' \, e^{-S[\phi'] + \int d^4x \, J(x)\phi'(x)} \, ,
\end{equation}
as the integration field \(\phi\) is merely a ``mute'' (dummy) integration variable; because it is functionally summed over all possible configurations, it can be freely renamed as \(\phi'\) without altering the value of the functional integral. This simple observation comes in handy, as we can define the renamed variable \(\phi'\) as a shifted or transformed version of the original field \(\phi\), providing the crucial starting point for all our subsequent derivations.

\subsubsection{Quantum Equations of Motion and the Ehrenfest Theorem}

Now, let us consider an infinitesimal local shift of the integration variable by taking:
\begin{equation} \label{eq:infinitesimal_field_shift}
    \phi'(x) = \phi(x) + \epsilon(x) \, ,
\end{equation}
with \(\epsilon(x)\) being an arbitrary, infinitesimal spacetime function. Assuming that the path integral measure is translationally invariant in field space (\(\mathcal{D}\phi' = \mathcal{D}\phi\)), we substitute \eqref{eq:infinitesimal_field_shift} into the right-hand side of \eqref{eq:path_integral_mute_variable}. Expanding the integrand to lowest (linear) order in \(\epsilon(x)\)---which is then stripped off at the end, say by taking the functional derivative \(\frac{\delta}{\delta \epsilon(x)}\)---we find the master identity:
\begin{equation} \label{eq:master_dyson_identity}
    \int \mathcal{D}\phi \left( -\frac{\delta S[\phi]}{\delta \phi(x)} + J(x) \right) e^{-S[\phi] + \int d^4y \, J(y)\phi(y)} = 0 \, .
\end{equation}
Dividing this expression by the vacuum normalization \(Z[0]\) to obtain properly normalized correlation functions, and subsequently setting the external source to vanish, \(J(x) = 0\), we immediately obtain:
\begin{equation} \label{eq:quantum_ehrenfest}
    \left\langle \frac{\delta S[\phi]}{\delta \phi(x)} \right\rangle = 0 \, ,
\end{equation}
indicating that the classical equations of motion hold exactly as a quantum expectation value. Here, we have returned to denoting by \(\langle \cdots \rangle\) the fully normalized correlation function. Physically, this result represents the field-theoretic generalization of the \textbf{Ehrenfest theorem} from ordinary quantum mechanics (which states that the expectation values of the position and momentum operators obey the classical Hamiltonian equations of motion): in quantum field theory, the classical equations of motion are satisfied in expectation value by the quantum field operator.

\subsubsection{Schwinger--Dyson Equations}

Rather than setting \(J(x) = 0\) immediately in \eqref{eq:master_dyson_identity}, one can first take \(n\) functional derivatives with respect to the source, \(\frac{\delta}{\delta J(x_1)} \cdots \frac{\delta}{\delta J(x_n)}\). These functional differential equations for the generating functional \(Z[J]\), after setting the source to vanish at the end, yield the celebrated \textbf{Schwinger--Dyson equations}, which form an infinite tower of exact quantum relations connecting different correlation functions to one another:
\begin{equation} \label{eq:schwinger_dyson}
    \boxed{\left\langle \frac{\delta S[\phi]}{\delta \phi(x)} \phi(x_1) \cdots \phi(x_n) \right\rangle = \sum_{k=1}^n \left\langle \phi(x_1) \cdots \delta^4(x - x_k) \cdots \phi(x_n) \right\rangle} \, ,
\end{equation}
where on the right-hand side the Dirac delta function \(\delta^4(x - x_k)\) explicitly substitutes the field \(\phi(x_k)\) at the \(k\)-th position. To understand the origin of this contact term, note that in \eqref{eq:master_dyson_identity} the source \(J(x)\) multiplies the exponential \(e^{\int J \phi}\); when the functional derivative \(\frac{\delta}{\delta J(x_k)}\) acts directly on the prefactor \(J(x)\) rather than on the exponential, it produces \(\frac{\delta J(x)}{\delta J(x_k)} = \delta^4(x - x_k)\) without bringing down a factor of \(\phi(x_k)\) from the exponent. Therefore, \(\phi(x_k)\) does not appear in the final expression of that specific term, as the delta function sets \(x_k\) directly to the spacetime point \(x\) where the equation of motion operator is evaluated. These Schwinger--Dyson equations are non-perturbative identities that must be satisfied by all correlation functions of the theory.

As a concrete check, let us consider the free Euclidean Klein--Gordon theory with action \(S[\phi] = \int d^4x \, \frac{1}{2}\phi(-\partial^2 + m^2)\phi\). Evaluating \eqref{eq:schwinger_dyson} for \(n=1\) leads to:
\begin{equation} \label{eq:klein_gordon_sd}
    \left\langle \frac{\delta S[\phi]}{\delta \phi(x)} \phi(y) \right\rangle \equiv (-\partial_x^2 + m^2) \langle \phi(x)\phi(y) \rangle = \delta^4(x - y) \, ,
\end{equation}
where we recognize that the two-point correlation function satisfies the defining equation for the Green's function (the inverse) of the kinetic operator \(-\partial_x^2 + m^2\), thereby identifying the two-point function as the free propagator of the theory.

\subsection{Symmetries and their Consequences}

Let us now systematically study the quantum consequences of symmetries within the path integral. Suppose we have a finite transformation of the fields:
\begin{equation} \label{eq:finite_symmetry_transform}
    \phi(x) \rightarrow \phi'(x) = G(\phi, \partial_\mu\phi, x) \, ,
\end{equation}
such that the classical action is invariant under this mapping:
\begin{equation} \label{eq:action_invariance_finite}
    S[\phi'] = S[\phi] \, .
\end{equation}

We compute the generating functional \(Z[J]\) through the following chain of equalities:
\begin{equation} \label{eq:z_symmetry_derivation}
    Z[J] = \int \mathcal{D}\phi \, e^{-S[\phi] + \int d^4x \, J(x)\phi(x)} = \int \mathcal{D}\phi' \, e^{-S[\phi'] + \int d^4x \, J(x)\phi'(x)} = \int \mathcal{D}\phi \, e^{-S[\phi] + \int d^4x \, J(x)\phi'(x)} \, .
\end{equation}
Let us carefully trace the logic of these steps: in \(Z[J]\), we first renamed the dummy integration field \(\phi\) to be the transformed field \(\phi'\) under the symmetry transformation. Then, in the final step, we used the fact that the classical action is invariant under the symmetry transformation (\(S[\phi'] = S[\phi]\)) and assumed that the path integral measure is also invariant (\(\mathcal{D}\phi' = \mathcal{D}\phi\)) to obtain the final expression (if this last step fails because the quantum measure is not invariant under a classical symmetry, one encounters quantum \textbf{anomalies}).

By taking \(n\) functional derivatives of these path-integral identities with respect to the source \(J(x)\) and setting \(J=0\), one finds that the quantum correlation functions of the theory are invariant under the symmetry transformation:
\begin{equation} \label{eq:correlator_invariance_finite}
    \boxed{\langle \phi'(x_1) \cdots \phi'(x_n) \rangle = \langle \phi(x_1) \cdots \phi(x_n) \rangle} \, .
\end{equation}

\subsubsection{Infinitesimal Symmetries and Correlation Functions}

Let us now specialize our discussion to the case of infinitesimal continuous symmetries:
\begin{equation} \label{eq:infinitesimal_symmetry_def}
    \delta_\alpha \phi(x) = \alpha F(\phi, \partial_\mu\phi, x) \, ,
\end{equation}
where \(\alpha\) is an infinitesimal constant parameter. The classical action is invariant, meaning its first variation vanishes:
\begin{equation} \label{eq:infinitesimal_action_invariance}
    \delta_\alpha S[\phi] = 0 \, .
\end{equation}
Repeating the previous path-integral steps with the infinitesimal field redefinition:
\begin{equation} \label{eq:infinitesimal_field_redef}
    \phi'(x) = \phi(x) + \delta_\alpha \phi(x) \, ,
\end{equation}
and expanding the rightmost side of \eqref{eq:z_symmetry_derivation} to linear order in \(\alpha\), we find that the zeroth-order terms cancel \(Z[J]\), leaving:
\begin{equation} \label{eq:linear_order_alpha_integral}
    \int \mathcal{D}\phi \left( \int d^4x \, J(x) \delta_\alpha \phi(x) \right) e^{-S[\phi] + \int d^4y \, J(y)\phi(y)} = 0 \, ,
\end{equation}
which can be written in the equivalent, compact expectation-value form in the presence of the source \(J\):
\begin{equation} \label{eq:source_ward_identity}
    \left\langle \int d^4x \, J(x) \delta_\alpha \phi(x) \right\rangle_J = 0 \, .
\end{equation}
Taking \(n\) functional derivatives with respect to \(J(x_1), \dots, J(x_n)\) and setting the source to vanish at the end, we obtain the fundamental symmetry relation for correlation functions:
\begin{equation} \label{eq:infinitesimal_correlator_invariance}
    \boxed{\sum_{k=1}^n \left\langle \phi(x_1) \cdots \delta_\alpha \phi(x_k) \cdots \phi(x_n) \right\rangle = 0} \, .
\end{equation}
This relation is immediately recognized as the infinitesimal version of \eqref{eq:correlator_invariance_finite}, obtained by expanding \(\langle (\phi + \delta_\alpha \phi) \cdots (\phi + \delta_\alpha \phi) \rangle = \langle \phi \cdots \phi \rangle\) to first order in \(\alpha\).

\subsection{Ward Identities and Noether Current}

Finally, let us consider the central role of the conserved Noether current in deriving local Ward identities. As is well known from classical field theory (via the Noether trick), the Noether current can be systematically obtained by extending the global infinitesimal symmetry transformation of \eqref{eq:infinitesimal_symmetry_def} to a more general, local transformation:
\begin{equation} \label{eq:local_noether_transform}
    \delta_{\alpha(x)} \phi(x) = \alpha(x) F(\phi, \partial_\mu\phi, x) \, ,
\end{equation}
that contains an arbitrary smooth spacetime function \(\alpha(x)\). In general, because \(\alpha(x)\) is no longer constant, this local transformation is no longer a symmetry of the classical action. However, its variation must be proportional to the spacetime gradient \(\partial_\mu \alpha(x)\), so the action transforms as:
\begin{equation} \label{eq:action_variation_noether}
    \delta_{\alpha(x)} S[\phi] = \int d^4x \, \partial_\mu \alpha(x) \, j^\mu(x) \, ,
\end{equation}
since for a constant parameter \(\alpha\) (\(\partial_\mu \alpha = 0\)) one must recover the exact global symmetry, i.e., \eqref{eq:infinitesimal_action_invariance}. The vector \(j^\mu(x)\) is identified as the classical Noether current.

Now, running again through the previous set of path-integral steps using the local change of integration variables:
\begin{equation} \label{eq:local_field_redef}
    \phi'(x) = \phi(x) + \delta_{\alpha(x)} \phi(x) \, ,
\end{equation}
we find (assuming again the invariance of the functional measure, \(\mathcal{D}\phi' = \mathcal{D}\phi\)):
\begin{equation} \label{eq:z_local_shift_1}
    Z[J] = \int \mathcal{D}\phi \, e^{-S[\phi] + \int d^4x \, J(x)\phi(x)} = \int \mathcal{D}\phi' \, e^{-S[\phi'] + \int d^4x \, J(x)\phi'(x)}
\end{equation}
\begin{equation} \label{eq:z_local_shift_2}
    = \int \mathcal{D}\phi \, e^{-S[\phi] - \delta_{\alpha(x)}S[\phi] + \int d^4x \, J(x)\phi(x) + \int d^4x \, J(x)\delta_{\alpha(x)}\phi(x)} \, ,
\end{equation}
which, expanding the exponential to linear order in the arbitrary function \(\alpha(x)\), implies:
\begin{equation} \label{eq:local_ward_integral}
    \int \mathcal{D}\phi \left( \delta_{\alpha(x)}S[\phi] - \int d^4x \, J(x) \delta_{\alpha(x)}\phi(x) \right) e^{-S[\phi] + \int d^4y \, J(y)\phi(y)} = 0 \, ,
\end{equation}
i.e., written in terms of the sourced expectation value:
\begin{equation} \label{eq:local_ward_expectation}
    \left\langle \delta_{\alpha(x)}S[\phi] - \int d^4x \, J(x) \delta_{\alpha(x)}\phi(x) \right\rangle_J = 0 \, .
\end{equation}

Substituting \eqref{eq:action_variation_noether} and integrating by parts to move the derivative \(\partial_\mu\) onto the current \(j^\mu(x)\), one can remove the arbitrary function \(\alpha(x)\) (say, by taking a functional derivative \(\frac{\delta}{\delta \alpha(x)}\)) and write this relation in a purely local form as:
\begin{equation} \label{eq:local_ward_with_source}
    \partial_\mu \langle j^\mu(x) \rangle_J = J(x) \langle \bar{\delta}\phi(x) \rangle_J \, ,
\end{equation}
where we have defined the stripped variation \(\bar{\delta}\phi(x) = F(\phi, \partial_\mu\phi, x)\), obtained by setting \(\delta_{\alpha(x)}\phi(x) \equiv \alpha(x)\bar{\delta}\phi(x)\) to factor out the function \(\alpha(x)\) from \eqref{eq:local_noether_transform}. Taking \(n\) further functional derivatives \(\frac{\delta}{\delta J(x_1)} \cdots \frac{\delta}{\delta J(x_n)}\) of \eqref{eq:local_ward_with_source}, and then setting the source \(J\) to vanish, one gets the celebrated \textbf{Ward--Takahashi identity}:
\begin{equation} \label{eq:ward_takahashi_identity}
    \boxed{\frac{\partial}{\partial x^\mu} \left\langle j^\mu(x) \phi(x_1) \cdots \phi(x_n) \right\rangle + \sum_{k=1}^n \delta^4(x - x_k) \left\langle \phi(x_1) \cdots \bar{\delta}\phi(x_k) \cdots \phi(x_n) \right\rangle = 0} \, .
\end{equation}
This is the fundamental Ward identity that describes the conservation of the Noether current in the quantum theory: the divergence of the Noether current inside a correlation function vanishes everywhere except at the contact points \(x = x_k\), where it generates the symmetry variations of the inserted fields. Integrating \eqref{eq:ward_takahashi_identity} over all spacetime \(\int d^4x\), and assuming that the boundary terms at infinity vanish, the divergence term drops out and one immediately recovers the global symmetry identity of \eqref{eq:infinitesimal_correlator_invariance}.

Further technical details on this derivation may be found in the textbook by M. Srednicki, ``Quantum Field Theory'', chapter 22.

\subsection{Symmetries and Generating Functionals}

A continuous symmetry of the action tightly constrains both the connected correlation functions generated by \(W[J]\) and the one-particle-irreducible (1PI) proper vertices generated by the effective action \(\Gamma[\varphi]\). We will now derive these functional constraints directly from the single change of variables in the path integral established above. As we shall see, the main structural distinction is between \textit{linear} symmetry transformations, which act on the classical mean fields \(\varphi\) in the exact same functional way as they act on the elementary quantum fields \(\phi\), and \textit{nonlinear} transformations (such as BRST transformations in non-abelian gauge theories), whose quantum expectation values inevitably contain non-trivial loop fluctuations.

Throughout this section, \(\langle \cdots \rangle_J\) denotes a quantum expectation value normalized by the sourced partition function \(Z[J]\), evaluated in the presence of a non-vanishing external source \(J\):
\[\boxed{\langle \mathcal{O}[\phi] \rangle_J = \frac{1}{Z[J]} \int \mathcal{D}\phi \, \mathcal{O}[\phi] \, e^{-S[\phi] + \int d^4y \, J_a(y)\phi^a(y)}} \, .\]
With this normalization convention, \(\langle 1 \rangle_J = 1\) identically, and the classical mean field in the Euclidean convention (\(Z[J] = e^{W[J]}\)) is given by:
\[\varphi^a(x) = \langle \phi^a(x) \rangle_J = \frac{\delta \ln Z[J]}{\delta J_a(x)} = \frac{\delta W[J]}{\delta J_a(x)} \, .\]

The basic global Ward identity derived previously was written in \eqref{eq:source_ward_identity}, which we now extend to a general multicomponent set of bosonic fields \(\phi^a(x)\) possessing an infinitesimal global symmetry:
\begin{equation} \label{eq:multicomponent_global_sym}
    \delta_\alpha \phi^a(x) = \alpha F^a[\phi](x) \, , \qquad \delta_\alpha S = 0 \, ,
\end{equation}
leading at first order in the constant parameter \(\alpha\) to the fundamental functional constraint:
\begin{equation} \label{eq:functional_ward_constraint}
    \boxed{\int d^4x \, J_a(x) \left\langle F^a[\phi](x) \right\rangle_J = 0} \, .
\end{equation}
Notice the physical interpretation of this equation: the external source term \(\int J_a \phi^a\) generally breaks the symmetry of the full integrand in the path integral. Nevertheless, \eqref{eq:functional_ward_constraint} represents the exact quantum constraint that survives precisely because a mere change of dummy integration variables cannot change the value of the functional integral.

\subsubsection{Linear Symmetries: From \(W\) to \(\Gamma\)}

We now consider linear symmetries, taking \(F^a[\phi](x) = T^a_{\phantom{a}b}\phi^b(x)\), with a constant, field-independent generator matrix \(T\). Linearity allows us to take the quantum expectation value immediately, yielding \(\langle F^a \rangle_J = T^a_{\phantom{a}b}\varphi^b\).

Let us derive the resulting functional identity explicitly from the path integral. Keep the external source \(J\) fixed and change the integration variables according to the infinitesimal transformation:
\begin{equation} \label{eq:linear_field_change}
    \phi^{\prime a}(x) = \phi^a(x) + \alpha T^a_{\phantom{a}b}\phi^b(x) \, .
\end{equation}

Using the invariance of the classical action and the functional integration measure, we expand to first order in the constant parameter \(\alpha\):
\begin{equation} \label{eq:z_linear_expansion}
    \begin{aligned}
        Z[J] & = \int \mathcal{D}\phi' \, e^{-S[\phi'] + \int d^4x \, J_a \phi^{\prime a}}                                                                                             \\
             & = \int \mathcal{D}\phi \, e^{-S[\phi] + \int d^4x \, J_a \phi^a} \left[ 1 + \alpha \int d^4y \, J_a(y) T^a_{\phantom{a}b}\phi^b(y) \right] + \mathcal{O}(\alpha^2) \, .
    \end{aligned}
\end{equation}

The first equality simply renames the dummy integration variable. In the second, we express \(\phi'\) in terms of \(\phi\); only the source coupling term varies at first order in \(\alpha\). Subtracting \(Z[J]\) from both sides and setting the coefficient of \(\alpha\) to zero gives:
\begin{equation} \label{eq:linear_z_constraint_integral}
    0 = \int d^4y \, J_a(y) T^a_{\phantom{a}b} \int \mathcal{D}\phi \, \phi^b(y) \, e^{-S[\phi] + \int d^4x \, J_a \phi^a} \, .
\end{equation}

An insertion of the operator \(\phi^b(y)\) inside the path integral is generated by taking a functional derivative with respect to its source component \(J_b(y)\):
\begin{equation} \label{eq:source_derivative_phi}
    \frac{\delta Z[J]}{\delta J_b(y)} = \int \mathcal{D}\phi \, \phi^b(y) \, e^{-S[\phi] + \int d^4x \, J_a \phi^a} \, .
\end{equation}

Therefore, \eqref{eq:linear_z_constraint_integral} can be written as a differential constraint on the partition function:
\begin{equation} \label{eq:linear_z_differential}
    \int d^4y \, J_a(y) T^a_{\phantom{a}b} \frac{\delta Z[J]}{\delta J_b(y)} = 0 \, .
\end{equation}

Finally, dividing by \(Z[J]\) and using the chain rule relation for the connected generator \(W[J] = -i \ln Z[J]\), namely:
\begin{equation} \label{eq:z_to_w_derivative_relation}
    \frac{1}{Z[J]} \frac{\delta Z[J]}{\delta J_b(y)} = \frac{\delta \ln Z[J]}{\delta J_b(y)} = \frac{\delta W[J]}{\delta J_b(y)} \, ,
\end{equation}
and renaming the dummy integration variable \(y\) as \(x\), we obtain the fundamental constraint satisfied by the connected functional \(W[J]\):
\begin{equation} \label{eq:linear_ward_W}
    \boxed{\int d^4x \, J_a(x) T^a_{\phantom{a}b} \frac{\delta W}{\delta J_b(x)} = 0} \, .
\end{equation}
This exact same identity could have also been obtained directly by substituting the expectation value \(\langle F^a \rangle_J = T^a_{\phantom{a}b}\varphi^b\) into the general functional Ward constraint of \eqref{eq:functional_ward_constraint}.

There is an equivalent and illuminating interpretation in terms of transforming the source itself. If \(\phi' = M\phi\), where \(M = 1 + \alpha T + \mathcal{O}(\alpha^2)\), we can assign to the source a dual transformation, \(J' = M^{-1, T} J\). Then the source-field pairing is invariant, \(J'_a \phi^{\prime a} = J_a \phi^a\), so the change of variables leaves both the partition function and the connected generator invariant, \(Z[J'] = Z[J]\) and \(W[J'] = W[J]\). Infinitesimally, this implies the source variation:
\begin{equation} \label{eq:source_infinitesimal_transform}
    \delta_\alpha J_a = -\alpha J_b T^b_{\phantom{b}a} \, .
\end{equation}

Differentiating the invariance relation \(W[J'] = W[J]\) shows that the corresponding mean fields obey \(\varphi[J'] = M\varphi[J]\). Since the Legendre transform pairs \(J'\phi' = J\phi\), the quantum effective action \(\Gamma[\varphi]\) is invariant under the linear transformation as well. More directly, substituting the Legendre conjugate relation \(\frac{\delta \Gamma}{\delta \varphi^a(x)} = -J_a(x)\) (from \eqref{eq:basic_legendre_properties}) into \eqref{eq:linear_ward_W} gives:
\begin{equation} \label{eq:linear_ward_Gamma}
    \boxed{\int d^4x \, \frac{\delta \Gamma}{\delta \varphi^a(x)} T^a_{\phantom{a}b} \varphi^b(x) = 0} \, .
\end{equation}
This is precisely the statement that \(\delta_\alpha \Gamma = 0\) for the infinitesimal field variation \(\delta_\alpha \varphi^a = \alpha T^a_{\phantom{a}b}\varphi^b\). Thus, an anomaly-free linear symmetry of the classical action is strictly preserved as a symmetry of the full quantum effective action, including all quantum loop corrections.

To see explicitly what these functional identities imply for physical correlation functions, let us define the connected two-point function (the exact connected propagator \(G_c^{ab}(x,y; J)\), already described in Minkowski time in \eqref{eq:exact_propagator_inversion}):
\begin{equation} \label{eq:connected_two_point_def}
    G_c^{ab}(x, y; J) = \frac{\delta^2 W}{\delta J_a(x) \delta J_b(y)} = \langle \phi^a(x)\phi^b(y) \rangle_J - \varphi^a(x)\varphi^b(y) \, .
\end{equation}
Differentiating the connected Ward identity of \eqref{eq:linear_ward_W} twice with respect to the sources and then setting the external source \(J = 0\) gives:
\begin{equation} \label{eq:two_point_ward_identity}
    T^a_{\phantom{a}c} G_c^{cb}(x, y; 0) + T^b_{\phantom{a}c} G_c^{ac}(x, y; 0) = 0 \, .
\end{equation}
Taking further functional derivatives generates analogous symmetry rules for all higher-order connected correlation functions.

Similarly, derivatives of the effective-action identity in \eqref{eq:linear_ward_Gamma} constrain the 1PI proper vertices. For clarity, let the index \(i\) comprehensively denote both a discrete field index and a continuous spacetime coordinate, so that a repeated index implies implicit summation and spacetime integration. Writing \(\Gamma_{,i} \equiv \frac{\delta \Gamma}{\delta \varphi^i}\), taking two functional derivatives evaluated at a constant background field configuration \(v\) yields the relation:
\begin{equation} \label{eq:three_point_vertex_ward}
    \Gamma^{(3)}_{ijk}(v) (T v)^i + \Gamma^{(2)}_{ij}(v) T^i_{\phantom{i}k} + \Gamma^{(2)}_{ik}(v) T^i_{\phantom{i}j} = 0 \, .
\end{equation}
At a true symmetry-preserving background where \(Tv = 0\) (such as \(v = 0\) in an unbroken phase), the first term vanishes and the exact two-point vertex (the inverse propagator) is forced to satisfy a commutation relation with the generator matrix, meaning it forms an invariant tensor under the symmetry group. At a noninvariant background, the additional linear term must be systematically retained.

The effective-action identity extends directly to an affine symmetry of the form \(F^a[\phi](x) = T^a_{\phantom{a}b}\phi^b(x) + f^a(x)\), where the inhomogeneous shift function \(f^a(x)\) is completely field-independent:
\begin{equation} \label{eq:affine_ward_Gamma}
    \int d^4x \, \frac{\delta \Gamma}{\delta \varphi^a(x)} \Big( T^a_{\phantom{a}b}\varphi^b(x) + f^a(x) \Big) = 0 \, .
\end{equation}
The underlying reason is again that the linear quantum expectation value obeys \(\langle F^a[\phi] \rangle_J = F^a[\varphi]\).

\subsubsection{Nonlinear Symmetries: Why Fluctuations Enter}

For a genuinely nonlinear transformation, taking a naive quantum expectation value does not commute with evaluating the functional \(F\) at the quantum operator level because of quantum fluctuations. For example, consider a quadratic term in a symmetry variation, \(F^a[\phi](x) = C^a_{\phantom{a}bc}\phi^b(x)\phi^c(x)\). Its exact expectation value in the presence of sources evaluates to:
\begin{equation} \label{eq:nonlinear_expectation_value}
    \langle F^a[\phi](x) \rangle_J = C^a_{\phantom{a}bc} \left[ \varphi^b(x)\varphi^c(x) + G_c^{bc}(x, x; J) \right] \, .
\end{equation}
The crucial extra term appearing here is precisely the connected two-point function evaluated at coincident points (the equal-time loop fluctuation or propagator contraction about the mean field).

A vital technical point must be emphasized here: coincident operator products in quantum field theory are divergent and must be strictly regulated and renormalized consistently as composite operators.

Defining the exact quantum-induced functional as the full expectation value expressed in terms of the mean field \(\varphi\):
\begin{equation} \label{eq:induced_functional_F}
    \mathcal{F}^a[\varphi](x) \equiv \langle F^a[\phi](x) \rangle_{J(\varphi)} \, ,
\end{equation}
the general effective-action identity for nonlinear symmetries takes the compact form:
\begin{equation} \label{eq:nonlinear_ward_Gamma}
    \boxed{\int d^4x \, \frac{\delta \Gamma}{\delta \varphi^a(x)} \mathcal{F}^a[\varphi](x) = 0} \, .
\end{equation}
In other words, the effective action \(\Gamma\) is invariant under the induced infinitesimal field variation \(\delta_\alpha \varphi^a = \alpha \mathcal{F}^a[\varphi]\), but this does \textit{not} in general imply simple invariance under the naive classical transformation rule \(\delta_\alpha \phi^a = \alpha F^a[\phi]\) due to the presence of quantum fluctuation loop terms.

\subsection{Sources for Composite Operators}

The functional identities derived above already successfully express the exact symmetry constraints on quantum correlation functions. However, introducing auxiliary external sources coupled directly to composite operators provides a powerful bookkeeping device when one wishes to treat composite insertions like \(\mathcal{F}\) as systematic functional derivatives.

To implement this, we generalize the generating functional by adding a source term coupled to the composite symmetry variation:
\begin{equation} \label{eq:composite_source_Z_W}
    Z[J, K] = \int \mathcal{D}\phi \, e^{-S[\phi] + \int d^4x \, (J_a \phi^a + K_a F^a)} \, , \qquad W[J, K] = \ln Z[J, K] \, ,
\end{equation}
where the new source \(K_a(x)\) couples directly to the nonlinear symmetry variation \(F^a[\phi]\). We then define the effective action via a Legendre transform with respect to the primary source \(J_a\) only:
\begin{equation} \label{eq:composite_source_Gamma}
    \Gamma[\varphi, K] = \int d^4x \, J_a \varphi^a - W[J, K] \, , \qquad \varphi^a(x) = \frac{\delta W}{\delta J_a(x)} \, .
\end{equation}
In this formulation, the auxiliary source \(K_a(x)\) remains treated as an independent background argument throughout the Legendre transform. At fixed background field \(\varphi\), terms involving the implicit variation of \(J\) with respect to \(K\) cancel out, yielding the fundamental integro-differential relation:
\begin{equation} \label{eq:gamma_K_derivative_relation}
    \left. \frac{\delta \Gamma}{\delta K_a(x)} \right|_{\varphi = \text{const}} = - \int d^4y \, \frac{\delta W}{\delta K_a(x)} \frac{\delta K_a}{\delta K} \dots = -\left\langle F^a[\phi](x) \right\rangle_{J,K} = -\mathcal{F}^a[\varphi](x) \, .
\end{equation}
Consequently, evaluating the functional derivative of the effective action with respect to the source \(K_a\) at the point where \(K = 0\) directly generates the exact composite operator insertion in terms of the mean field \(\varphi\):
\begin{equation} \label{eq:composite_source_master_identity}
    \left. \frac{\delta \Gamma}{\delta K_a(x)} \right|_{K=0} = -\mathcal{F}^a[\varphi](x) \, .
\end{equation}

Substituting this back into the generalized nonlinear invariance condition yields the master equation:
\begin{equation} \label{eq:master_composite_ward}
    \int d^4x \, \frac{\delta \Gamma}{\delta \varphi^a(x)} \left. \frac{\delta \Gamma}{\delta K_a(x)} \right|_{K=0} = 0 \, .
\end{equation}

These represent alternative functional expressions of the exact quantum symmetry constraints. They are formally asserted here at \(K = 0\); at nonzero composite sources, the transformation of variables also induces non-trivial variations proportional to \(K_a F^a\). In a fully renormalized quantum field theory, these sources and composite insertions require a consistent composite-operator renormalization prescription, including all necessary local counterterms.

For the profound case of BRST symmetry (which is a nilpotent, nonlinear, gauge-fixing dependent symmetry), the analogous functional construction utilizes Grassmann-graded fields and specifies left/right functional derivations, ultimately producing the celebrated \textbf{Slavnov--Taylor identities}; note that the bosonic functional formulas derived above do not explicitly capture the sign structures associated with fermionic ghost fields.